\documentclass[10.5pt,a4paper]{article}
\usepackage[utf8]{inputenc}
\usepackage[T1]{fontenc}
\usepackage[french]{babel}
\usepackage{amsmath,amssymb}
\usepackage[margin=1.9cm]{geometry}
\usepackage{enumitem}
\setlist{itemsep=2pt,topsep=3pt}
\setlength{\parindent}{0pt}
\setlength{\parskip}{4pt}

\begin{document}

\begin{center}
{\Large\bfseries Fiche de lecture -- Projet 9 : Deep Hedging}\\[2pt]
Module Finance des Marchés et ERP -- Groupe 9\\
Membres : Marouane, Aymane, Sara, Meryem
\end{center}

\section*{1. Formules clés}

\textbf{Dynamique du sous-jacent (GBM, marché simulé).}
\[
S_{k+1} = S_k \exp\!\Big[(\mu-\tfrac{\sigma^2}{2})\Delta t + \sigma\sqrt{\Delta t}\,Z_k\Big],\qquad Z_k\sim\mathcal N(0,1),\ \Delta t=T/n
\]
Le prix évolue de façon aléatoire log-normale ; c'est le « marché » sur lequel on entraîne et on teste les stratégies.

\textbf{Payoff de l'option vendue.}
\[
Z = \max(S_T-K,\,0)\quad(\text{call européen})
\]
C'est ce que le vendeur doit payer à l'échéance : c'est la perte qu'on cherche à compenser avec la couverture.

\textbf{Delta de Black-Scholes.}
\[
\Delta_k = N(d_1),\qquad d_1=\frac{\ln(S_k/K)+(r+\sigma^2/2)(T-t_k)}{\sigma\sqrt{T-t_k}}
\]
Le Delta est le nombre d'actions à détenir à la date $t_k$ pour neutraliser la sensibilité de l'option au prix du sous-jacent ; c'est la stratégie classique de référence.

\textbf{P\&L de la couverture avec coûts de transaction.}
\[
\text{P\&L} = -Z + \sum_{k=0}^{n-1}\delta_k\,(S_{k+1}-S_k) - \sum_{k}c\,|\delta_k-\delta_{k-1}|\,S_k,\qquad \delta_{-1}=0
\]
On additionne le gain (ou la perte) du portefeuille d'actions, on retire le payoff dû et les coûts proportionnels $c$ payés à chaque réajustement de la position.

\textbf{Réseau de neurones.}
\[
\delta_k = f_\theta\big(\ln(S_k/K),\ T-t_k,\ \delta_{k-1}\big)
\]
Le réseau reçoit l'état du marché (prix relatif, temps restant, position précédente) et renvoie directement la quantité à détenir, sans formule analytique.

\textbf{Objectif d'entraînement.}
\[
\min_\theta\ \mathrm{Var}(\text{P\&L})\qquad(\text{ou } \mathrm{CVaR}_\alpha,\ \text{ou } \tfrac1\lambda\ln \mathbb E[e^{-\lambda\,\text{P\&L}}])
\]
On cherche les poids $\theta$ qui rendent le résultat final le plus stable possible : un P\&L peu dispersé autour de zéro signifie une bonne couverture.

\textbf{Hypothèses.} Marché simulé (GBM ou Heston) : le réseau n'est valable que dans la dynamique d'entraînement ; coûts $c$ connus ; fréquence de rééquilibrage fixe ; trajectoires d'entraînement indépendantes, avec séparation train/test pour mesurer la généralisation.

\section*{2. Lien avec le TP 0 (module Odoo \texttt{cours\_derives})}

On ne repart pas de zéro : on étend le module existant dans le même dossier d'addons.
\begin{itemize}
\item Nouveau modèle \texttt{cours.simulation.couverture} avec les champs : méthode de couverture (sélection « Classique Delta » / « IA »), paramètres de marché ($S_0$, $K$, $T$, $r$, $\sigma$), coût de transaction $c$, nombre de trajectoires ; et des champs résultats stockés (moyenne, écart-type, VaR du P\&L).
\item Bouton « Lancer la simulation » (méthode \texttt{action\_lancer\_simulation()}) qui appelle le code Python du dossier \texttt{src/} (module interne ou petit service API) et enregistre les résultats ; vue formulaire + vue graphique comparant les deux méthodes.
\item Le réseau est entraîné \emph{hors d'Odoo} (notebook / PyTorch) et sauvegardé dans un fichier \texttt{.pt} ; Odoo le charge et fait seulement de l'\textbf{inférence} (pas de ré-entraînement, trop lent).
\item On réutilise la fonction de prix et le Delta de Black-Scholes du Groupe 1 (\texttt{cours.option}) : s'accorder sur les noms de champs et les conventions (unités de $T$, annualisation).
\item Fichiers à mettre à jour : \texttt{models/}, \texttt{views/}, \texttt{security/ir.model.access.csv}, \texttt{\_\_manifest\_\_.py}.
\end{itemize}

\section*{3. Algorithme en Python}
\begin{enumerate}
\item \textbf{Simuler les prix} : générer $M$ trajectoires GBM vectorisées avec NumPy/PyTorch (ex. 100\,000 pour l'entraînement, 20\,000 pour le test, $n=30$ à $252$ pas), graine fixée pour la reproductibilité.
\item \textbf{Calculer le P\&L} : pour chaque date $k$, calculer $\delta_k$ (Delta de BS ou sortie du réseau), mettre à jour le gain sur les actions, retirer les coûts de transaction, puis soustraire le payoff à l'échéance.
\item \textbf{Entraîner le réseau} : MLP à 2 couches cachées (64 neurones, ReLU), optimiseur Adam (taux $10^{-3}$), perte = variance du P\&L, nouvelles trajectoires à chaque epoch, arrêt anticipé sur trajectoires de validation.
\item \textbf{Comparer les stratégies} : Delta de BS, aucune couverture / aléatoire, réseau ; indicateurs : moyenne, écart-type, quantile 5\,\%, coût de transaction total, histogrammes superposés.
\item \textbf{Analyser et intégrer} : tracer $\delta(S)$ du réseau contre le Delta de BS, étudier l'effet de $c$ (le réseau apprend à moins rééquilibrer quand les coûts augmentent), sauvegarder le modèle en \texttt{.pt} et l'appeler depuis Odoo.
\end{enumerate}

\end{document}
